\scanpage{17}
\mspage{15}

\deleted{tés, les morphismes \uline{toutes} les correspondances \inserted{co}homologiques (\unsure{e|o}n ve
du thorème de Thom rappelé plus haut). On trouve donc un foncteur
pleinement fidèle de cette catégorie dans \unsure{c|o}elle des vectoriels
gradués d\deleted{\unsure{e}} d positivement de dim. finie sur $\underline{F}_2$ , qui est déjà abéli
ne, et si on veut rendre comme plus haut $\underline{F}_2(-1)$ inversible, on trou
tous les vectoriels gradués de dim. finie sur $\underline{F}_2$. C'est un peu tr\unsure{o}
trivial comme structure, et on peut faire mieux,} \deleted{\unsure{e|E}n prenant comme}
\inserted{Pour trouver une catégorie universelle,}
\deleted{correspondances Pour c\unsure{e}ci} on associe à toute X
\deleted{l\unsure{e|a} \unsure{sous} \unsure{anneau} ensem\unsure{b}}
\deleted{\unsure{d|l}} l'ensemble \inserted{\begin{tabular}[b]{@{}l@{}}S(X)\\[-.6ex]\begin{tikzpicture}[x=1cm,y=1cm,baseline=0pt]\draw[line width=.4pt] (0,0)--(.03,.05);\draw[line width=.4pt] (.08,.13)--(.15,.2);\draw[line width=.4pt] (.22,.28)--(.43,.51);\draw[line width=.4pt] (.49,.58)--(.63,.75);\draw[line width=.4pt] (.67,.8)--(.69,.83);\end{tikzpicture}\end{tabular}} des sous-variétés \deleted{S(X)} de X, modulo équivalence hom\unsure{o}
gique. Donc S(X) $\subset$ H$^{\bullet}$(X), et on prouve qu'on trouve en fait un sous
anneau de H$^{\bullet}$(X) (grâce à un théorème de Thom qui caractérise S(X) \handwritten{e\unsure{n|t}}
termes d'opérations de Steenrod, cf plus bas). On prouve de même
(par le même théorème) que les S(XxY) sont stables par composition
correspondances cohomologiques( cela ne semble pas trivial, puisque
les S(X) ne sont pas stables par Gysin...). D'où une catégorie ad\unsure{d|c}
tive dont les objets sont ceux de C, les flêches sont les \inserted{\begin{tabular}[b]{@{}l@{}}\handwritten{él. des}\\[-.6ex]\begin{tikzpicture}[x=1cm,y=1cm,baseline=0pt]\draw[line width=.4pt] (0,.5) .. controls (.22,.28) and (.67,.32) .. (.99,.33) .. controls (1.16,.4) and (1.3,.44) .. (1.44,.37) .. controls (1.48,.32) and (1.47,.14) .. (1.53,.13) .. controls (1.66,.35) and (1.88,.49) .. (2.13,.48) .. controls (2.38,.42) and (2.61,.4) .. (2.77,.55) .. controls (2.86,.66) and (2.91,.76) .. (2.9,.84)--(2.78,.85);\draw[line width=.4pt] (1.53,.13)--(1.5,-.15);\end{tikzpicture}\end{tabular}} S(XxY),
Il semble qu'elle doive jouer le rôle universel que j'ai envisagé
plus haut. En tous cas, sa définition géo\unsure{\deleted{é}m|\deleted{m}é}\unsure{\deleted{m}é|\deleted{é}m}trique est fort simple
Nous allons en donner une interpétation algébrique. P\unsure{\u{o}}ur ce\unsure{\deleted{i}c|\deleted{c}i}\unsure{\deleted{c}i|\deleted{i}c},
considérons l'anneau des opérations cohomologiques \inserted{``stables''} \deleted{additives sur \unsure{l|i}}
\margin{\begin{tikzpicture}[x=1cm,y=1cm,baseline=0pt,line cap=round]\draw[line width=.4pt] (.27,4.8)--(.18,4.57);\draw[line width=.4pt] (.13,4.45)--(.12,4.4);\draw[line width=.4pt] (.08,4.34)--(.08,4.32);\draw[line width=.4pt] (.04,4.25)--(.03,4.2);\draw[line width=.4pt] (-.04,3.96)--(-.04,3.93);\draw[line width=.4pt] (-.05,3.82)--(-.07,3.66);\draw[line width=.4pt] (-.07,3.38)--(-.07,3.29);\draw[line width=.4pt] (-.07,3.2)--(-.07,3.14);\draw[line width=.4pt] (-.06,3.05)--(-.05,2.99);\draw[line width=.4pt] (-.01,2.75)--(0,2.67);\draw[line width=.4pt] (.01,2.57) .. controls (.06,2.36) and (.1,2.11) .. (.14,1.98);\draw[line width=.4pt] (.16,1.85)--(.17,1.65);\draw[line width=.4pt] (.18,1.5)--(.18,1.41);\draw[line width=.4pt] (.17,1.27)--(.17,1.18);\draw[line width=.4pt] (.16,1.02)--(.16,1);\draw[line width=.5pt] (.15,.87)--(.15,.77);\draw[line width=.3pt] (.16,.53)--(.16,.51);\draw[line width=.3pt] (.17,.35)--(.17,.33);\draw[line width=.3pt] (.175,.12)--(.175,.1);\begin{scope}[yshift=-4.2cm]\draw[line width=.25pt] (.05,4.13)--(.05,4.11);\draw[line width=.2pt] (0,3.8)--(0,3.78);\draw[line width=.25pt] (.08,3.15)--(.08,3.13);\draw[line width=.25pt] (.09,2.96)--(.09,2.94);\draw[line width=.4pt] (.11,2.42)--(.11,2.37);\draw[line width=.4pt] (.115,2.31)--(.115,2.25);\draw[line width=.2pt] (.115,2.2)--(.115,2.18);\draw[line width=.2pt] (.12,2.07)--(.12,2.05);\draw[line width=.4pt] (.14,1.68)--(.14,1.63);\draw[line width=.25pt] (.15,1.46)--(.15,1.44);\draw[line width=.4pt] (.14,1.08)--(.14,.92);\draw[line width=.4pt] (.11,.81)--(.1,.7);\draw[line width=.4pt] (.09,.51) .. controls (.09,.39) and (.1,.33) .. (.14,.29);\draw[line width=.2pt] (.19,.15)--(.19,.13);\end{scope}\end{tikzpicture}}
\deleted{des endomorphismes du}
\margin{\figure{\begin{tikzpicture}[x=1cm,y=1cm,baseline=0pt,line cap=round]\draw[line width=.5pt] (1,3.6)--(.91,3.45);\draw[line width=.5pt] (1.35,3.3) .. controls (1.32,3.2) and (1.27,3.08) .. (1.22,3.0);\draw[line width=.5pt] (.8,3.22) .. controls (.65,3.11) and (.43,2.9) .. (.5,2.81) .. controls (.56,2.75) and (.64,2.79) .. (.68,2.85);\draw[line width=.5pt] (.9,3.1) .. controls (1.02,3.2) and (1.1,2.99) .. (1.1,2.83) .. controls (1.1,2.61) and (1.12,2.39) .. (1.06,2.24) .. controls (.94,2.14) and (.77,2.13) .. (.68,2.19)--(.66,2.27);\draw[line width=.5pt] (.99,2.16)--(1.11,2.03)--(1.01,1.97);\draw[line width=.65pt] (.87,2.66)--(.88,2.68);\draw[line width=.5pt] (1,1.92) .. controls (.9,1.75) and (.83,1.6) .. (.77,1.48) .. controls (.67,1.35) and (.57,1.16) .. (.63,1.14);\draw[line width=.5pt] (.75,1.5) .. controls (.91,1.45) and (1.04,1.32) .. (1.13,1.32) .. controls (1.23,1.32) and (1.34,1.36) .. (1.4,1.43);\draw[line width=.5pt] (1.08,1.33)--(1.1,1.15);\draw[line width=.7pt] (1,.55)--(1.02,.58);\end{tikzpicture}}}
\deleted{\deleted{le} foncteur $X\quad H^{\bullet}(X)$ qui sont additifs, et qui en degré zéro se
réduisent à un multiple de l'identité}\begin{tikzpicture}[x=1cm,y=1cm,baseline=0pt]\draw[line width=.4pt] (0,0)--(.18,-.02) .. controls (.23,-.03) and (.25,-.1) .. (.31,-.06) .. controls (.48,.18) and (.63,.51) .. (.78,.83) .. controls (.84,.96) and (.85,1.02) .. (.79,1.08);\end{tikzpicture}
\figure{\begin{tikzpicture}[x=1cm,y=1cm,baseline=0pt]\draw[line width=.4pt] (0,.38) .. controls (.02,.24) and (0,.15) .. (.1,.06)--(.08,0);\draw[line width=.4pt] (.1,.07)--(.36,.33);\draw[line width=.4pt] (.4,.36)--(.45,.39);\draw[line width=.4pt] (.67,.39) .. controls (.71,.2) and (.78,.09) .. (.87,.11) .. controls (1.08,.12) and (1.29,.36) .. (1.45,.35) .. controls (1.56,.32) and (1.56,.14) .. (1.75,.14) .. controls (1.96,.14) and (2.26,.33) .. (2.4,.32) .. controls (2.54,.3) and (2.66,.04) .. (2.82,.04) .. controls (3.14,.04) and (3.48,.11) .. (3.61,.08)--(3.8,.06);\end{tikzpicture}}
C'est un anneau gradué, avec
application diagonale anticommutative \inserted{\begin{tabular}[b]{@{}l@{}}i.e. commutative (car. 2 !)\\[-.9ex]
\begin{tikzpicture}[x=1cm,y=1cm,baseline=0pt]\draw[line width=.4pt] (0,0) .. controls (.13,.12) and (.48,.13) .. (1.16,.13)--(1.18,.12)--(1.18,.075);\draw[line width=.4pt] (1.25,.12)--(1.34,.115);\draw[line width=.4pt] (1.4,.11)--(1.49,.1);\draw[line width=.4pt] (1.58,.09)--(1.7,.08);\draw[line width=.4pt] (1.8,.08)--(1.84,.08);\draw[line width=.4pt] (3.78,.1) .. controls (3.96,.11) and (4.12,.11) .. (4.22,.19);\draw[line width=.4pt] (4.27,.23)--(4.31,.29);\draw[line width=.4pt] (4.38,.33)--(4.4,.38)--(4.33,.38);\end{tikzpicture}
\end{tabular}}, associative et unitaire. So\unsure{n|i}
dual gradué est une algèbre \inserted{\begin{tabular}[b]{@{}l@{}}graduée anti\unsure{o|c}ommutative \begin{tikzpicture}[x=1cm,y=1cm,baseline=0pt]\draw[line width=.35pt](0,-.05)--(.09,-.05);\end{tikzpicture} donc \begin{tikzpicture}[x=1cm,y=1cm,baseline=0pt]\draw[line width=.35pt](0,-.05)--(.09,-.05);\end{tikzpicture}\deleted{commutative};\\[-.9ex]
\begin{tikzpicture}[x=1cm,y=1cm,baseline=0pt,line cap=round]\draw[line width=.5pt] (0,0)--(.01,.01);\draw[line width=.4pt] (.25,.03)--(.26,.03);\draw[line width=.4pt] (.32,.04)--(.64,.06);\draw[line width=.4pt] (.71,.065)--(.78,.07);\draw[line width=.4pt] (.86,.075)--(.96,.09);\draw[line width=.4pt] (1.06,.1)--(1.1,.105);\draw[line width=.4pt] (1.18,.11)--(1.3,.12);\draw[line width=.4pt] (1.42,.13)--(1.53,.14);\draw[line width=.4pt] (1.6,.145)--(1.83,.14);\draw[line width=.4pt] (1.91,.14)--(2.49,.13);\draw[line width=.4pt] (2.58,.14)--(2.62,.14);\draw[line width=.4pt] (2.67,.14)--(2.68,.14);\end{tikzpicture}
\end{tabular}} commutative,avec application diagonale
associative et unitaire; il est réduit à $\underline{F}_2$ en degré zéro. Donc il
corespond à un schéma en groupes affine G sur $\underline{F}_2$ , qui est d'aille\unsure{u}
le foncteur des $\unsure{\otimes|\Omega|Q}$-automorphismes \inserted{\begin{tabular}[b]{@{}l@{}}additifs\\[-.6ex]\begin{tikzpicture}[x=1cm,y=1cm,baseline=0pt]\draw[line width=.4pt] (0,.47)--(.01,.42);\draw[line width=.4pt] (.02,.39)--(.07,.33);\draw[line width=.4pt] (.12,.29) .. controls (.25,.21) and (.34,.18) .. (.5,.18) .. controls (.72,.17) and (.93,.18) .. (1.09,.14) .. controls (1.37,.1) and (1.5,.1) .. (1.76,.1) .. controls (1.86,.1) and (1.96,.05) .. (2.03,.04)--(2.08,-.08)--(2.02,.11)--(2.12,.2) .. controls (2.26,.13) and (2.3,.16) .. (2.39,.16) .. controls (2.73,.17) and (3.08,.1) .. (3.35,.39);\end{tikzpicture}\end{tabular}} du foncteur H$^{\bullet}$ sur C \handwritten{qui op\uline{èr}ent
tri\unsure{via}le\unsure{men}t sur \deleted{\illegible{1}} $H^{\unsure{\deleted{2}1|\deleted{1}2|q|0}}({\unsure{S|\lambda|\Delta}}^{\unsure{1|r|3}})$. \deleted{\illegible{1}}}

\margin{16}

